\section{Limitations}\label{sec:limitations}
\textbf{Reference quality.} The annotations contain known errors, and decisions about span boundaries or shared prices involve judgement. Some annotations began as AI suggestions that were reviewed manually. There was no systematic independent second annotation of every page. The additional visual review used one AI reviewer and checked emitted records, so it cannot establish how many visible offers were missed.

\textbf{Data and uncertainty.} Evaluation covers 42 PENNY pages, and test pages had already been inspected during development. The results therefore remain exploratory. The pipeline requires a PDF text layer, and transfer to new weeks or retailers was not tested. The bootstrap resamples one saved output set per system. It does not measure uncertainty from annotation errors, retraining or repeated model requests.

\textbf{Comparability.} The original API comparison covered models available through Academic Cloud and was not an exhaustive benchmark. Astra used shared context and tools in Codex, so its run differs from separate API requests. The recognition models also differ in pretraining, which prevents isolating the benefit of layout or images. Earlier grouping experiments used automatic references and partly incomplete run records.

\textbf{Evaluation scope.} Record F1 covers product names and regular prices, excludes incomplete fragments and does not assess all offer conditions. Additional post-processing of the final records, such as standardising quantity formats or handling duplicate offers, was not evaluated systematically. Its benefit therefore remains unmeasured. Runtime measurements excluded PDF preparation and compared a local machine with an uncontrolled remote service, including network delays. Total operating costs were not measured.
